% =============================================================================
%  IIIT RAICHUR -- B.Tech PROJECT THESIS
%  Department of Computer Science and Engineering
%  Individual Report: Shaurya Yadav (CS22B1050)
%  Emphasis: Game Architecture, Quest System, LLM Integration
% =============================================================================

\documentclass[12pt, a4paper, oneside]{report}

% ---------------------- PACKAGES ---------------------------------------------
\usepackage[utf8]{inputenc}
\usepackage[T1]{fontenc}
\usepackage{lmodern}
\usepackage[margin=1in, top=1.2in, bottom=1.2in]{geometry}
\usepackage{graphicx}
\usepackage{setspace}
\usepackage{titlesec}
\usepackage{tocloft}
\usepackage{fancyhdr}
\usepackage{booktabs}
\usepackage{longtable}
\usepackage{array}
\usepackage{tabularx}
\usepackage{multirow}
\usepackage{caption}
\usepackage{subcaption}
\usepackage{amsmath, amssymb}
\usepackage{xcolor}
\usepackage[hidelinks, breaklinks=true]{hyperref}
\usepackage{cleveref}
\usepackage{url}
\usepackage{enumitem}
\usepackage{float}
\usepackage[numbers]{natbib}

% ---------------------- USER INPUTS ------------------------------------------
\newcommand{\thesisTitle}{Cooperative Non-Player Character Behavior via Multi-Agent Reinforcement Learning in a Role-Playing Game Environment}
\newcommand{\thesisTitleDisplay}{COOPERATIVE NON-PLAYER CHARACTER BEHAVIOR \\ VIA MULTI-AGENT REINFORCEMENT LEARNING \\ IN A ROLE-PLAYING GAME ENVIRONMENT}
\newcommand{\studentName}{Shaurya Yadav}
\newcommand{\rollNumber}{CS22B1050}
\newcommand{\supervisorName}{Dr.~Supervisor Name}
\newcommand{\departmentName}{Computer Science and Engineering}
\newcommand{\submissionMonthYear}{May 2026}
\newcommand{\submissionLocation}{Raichur -- 584135}
\newcommand{\degreeName}{Bachelor of Technology}
\newcommand{\logoPath}{images/iiitr_logo.png}

% ---------------------- FORMATTING -------------------------------------------
\onehalfspacing

\titleformat{\chapter}[display]
  {\normalfont\huge\bfseries}{\chaptertitlename\ \thechapter}{20pt}{\Huge}
\titlespacing*{\chapter}{0pt}{-30pt}{30pt}

\pagestyle{fancy}
\fancyhf{}
\fancyfoot[C]{\thepage}
\renewcommand{\headrulewidth}{0pt}
\renewcommand{\footrulewidth}{0pt}

\fancypagestyle{plain}{%
  \fancyhf{}%
  \fancyfoot[C]{\thepage}%
  \renewcommand{\headrulewidth}{0pt}%
  \renewcommand{\footrulewidth}{0pt}%
}

\renewcommand{\cftchapfont}{\bfseries}
\renewcommand{\cftchappagefont}{\bfseries}

\newcommand{\frontmatterheading}[1]{%
  \begin{center}\Large\bfseries #1\end{center}
  \vspace{0.5cm}
}

\setlength{\parskip}{0.5em}
\setlength{\parindent}{1.5em}

% =============================================================================
\begin{document}

% ---------------------- TITLE PAGE -------------------------------------------
\begin{titlepage}
    \centering
    \thispagestyle{empty}
    \vspace*{0.2cm}

    {\Large\bfseries \thesisTitleDisplay \par}
    \vspace{0.8cm}

    {\normalsize A Project Report Submitted \par}
    {\normalsize in Partial Fulfilment of the Requirements \par}
    {\normalsize for the Degree of \par}
    \vspace{0.5cm}

    {\large\bfseries \MakeUppercase{\degreeName} \par}
    \vspace{0.3cm}

    {\normalsize in \par}
    \vspace{0.2cm}

    {\large\bfseries \MakeUppercase{\departmentName} \par}
    \vspace{0.5cm}

    {\normalsize\itshape by \par}
    \vspace{0.3cm}

    {\normalsize\bfseries \studentName{} (Roll No.~\rollNumber) \par}
    \vspace{0.5cm}

    \includegraphics[width=1\textwidth]{\logoPath}\par
    \vspace{0.5cm}

    {\normalsize\itshape to \par}
    \vspace{0.2cm}

    {\normalsize\bfseries DEPARTMENT OF \MakeUppercase{\departmentName} \par}
    {\normalsize\bfseries INDIAN INSTITUTE OF INFORMATION TECHNOLOGY \par}
    {\normalsize\bfseries RAICHUR-584135, INDIA \par}
    \vspace{0.4cm}

    {\normalsize\itshape \submissionMonthYear \par}

\end{titlepage}

% ---------------------- FRONT MATTER (roman page numbers) --------------------
\pagenumbering{roman}
\setcounter{page}{2}

% ---------------------- DECLARATION ------------------------------------------
\clearpage
\thispagestyle{plain}
\frontmatterheading{DECLARATION}
\addcontentsline{toc}{chapter}{Declaration}

\noindent I, \textbf{\MakeUppercase{\studentName}} (Roll Number: \rollNumber), hereby
declare that this report entitled ``\textit{\thesisTitle}'' submitted to
Indian Institute of Information Technology Raichur towards partial requirement
of \degreeName{} in \departmentName{} is an original work carried out by me
under the supervision of Mentor, \supervisorName{}, and has not formed the
basis for the award of any degree or diploma, in this or any other institution
or university. I have sincerely tried to uphold the academic ethics and
honesty. Whenever an external information or statement or result is used then,
that has been duly acknowledged and cited.

\vspace{2.5cm}
\noindent
\begin{minipage}[t]{0.5\textwidth}
\submissionLocation \\
\submissionMonthYear
\end{minipage}
\hfill
\begin{minipage}[t]{0.4\textwidth}
\raggedleft
\studentName \\
Roll No.~\rollNumber
\end{minipage}

% ---------------------- CERTIFICATE ------------------------------------------
\clearpage
\thispagestyle{plain}
\frontmatterheading{CERTIFICATE}
\addcontentsline{toc}{chapter}{Certificate}

\noindent This is to certify that the work contained in this project report
entitled ``\textit{\thesisTitle}'' submitted by \textbf{\studentName} (Roll
No: \rollNumber) to the
Indian Institute of Information Technology Raichur towards partial requirement
of \degreeName{} in \departmentName{} has been carried out under my
supervision and that it has not been submitted elsewhere for the award of any
degree.

\vspace{3cm}
\noindent
\begin{minipage}[t]{0.5\textwidth}
\submissionLocation \\
\submissionMonthYear
\end{minipage}
\hfill
\begin{minipage}[t]{0.4\textwidth}
\raggedleft
(\supervisorName) \\
Project Supervisor
\end{minipage}

% ---------------------- ABSTRACT ---------------------------------------------
\clearpage
\thispagestyle{plain}
\frontmatterheading{ABSTRACT}
\addcontentsline{toc}{chapter}{Abstract}

This report presents the game architecture, quest system design, and LLM integration layer of a medieval village RPG research playground for studying emergent cooperative behavior in multi-agent systems. The game is structured around a hierarchical Markov Decision Process (MDP) with 7 quest stages and 16 static checkpoints, augmented by dynamic checkpoints that handle player deviations through nudging and convergence mechanisms. A 10-step turn processing pipeline orchestrates player actions, NPC decision-making via tabular Q-learning, environmental events, and telemetry logging within a coherent game loop. Six NPC agents develop adaptive personalities through a decomposed reward model $R_i(t) = P_i(t) + G_i(t) + \lambda_i(t) \cdot C(t)$ with a dynamic $\lambda$ feedback loop. An optional LLM integration layer, built as a thin HTTP adapter supporting Ollama, OpenRouter, and LM Studio, provides atmospheric narration across four use cases (input parsing, NPC dialogue, checkpoint generation, narration enhancement) with a 9-stage guardrail pipeline and deterministic template fallbacks ensuring full playability without model access. An ablation study across six conditions over 500 turns demonstrates that Q-learning drives cooperation from a 0.500 baseline to 0.851, the dynamic $\lambda$ feedback loop amplifies cooperation by 0.158, and LLM narration achieves higher contextual grounding (0.57 vs 0.49) and 2.6$\times$ richer vocabulary than static templates while maintaining output quality even with 4B-parameter models.

% ---------------------- TABLE OF CONTENTS ------------------------------------
\clearpage
\tableofcontents
\listoffigures
\addcontentsline{toc}{chapter}{List of Figures}
\listoftables
\addcontentsline{toc}{chapter}{List of Tables}

\clearpage
\pagenumbering{arabic}
\setcounter{page}{1}

% =============================================================================
% CHAPTER 1: INTRODUCTION
% =============================================================================
\chapter{Introduction}
\label{chap:introduction}

Non-player characters (NPCs) in modern games typically follow scripted behavior trees or finite state machines, producing predictable and repetitive interactions. This project explores whether reinforcement learning can produce emergent cooperative behavior in a multi-agent game environment without explicit scripting of cooperation.

We build a single-player medieval village RPG where six NPCs learn autonomously via tabular Q-learning. The central research question is:

\begin{quote}
\textit{Can NPCs learn to shift from individually optimal strategies to cooperative behavior when community welfare contributes to their total reward?}
\end{quote}

\section{Problem Statement}
\label{sec:problem-statement}

Current NPC systems in games rely on hand-authored behavior trees, finite state machines, or scripted dialogue. These approaches produce deterministic, repetitive behavior that breaks immersion and cannot adapt to player strategy. The core problems are:

\begin{itemize}
    \item \textbf{No emergent cooperation:} Scripted NPCs cannot discover cooperative strategies on their own --- any prosocial behavior must be explicitly programmed by designers for every scenario.
    \item \textbf{Brittleness to perturbation:} Static behavior trees cannot adapt when the game environment changes (e.g., resource shocks, new player strategies). NPCs continue executing their scripts regardless of context.
    \item \textbf{No individual--community trade-off:} Real social agents must balance personal gain against group welfare. Scripted NPCs have no mechanism to represent, let alone dynamically adjust, this trade-off.
    \item \textbf{Lack of interpretability in deep RL alternatives:} While deep RL can produce adaptive NPC behavior, the learned policies are opaque. For a research playground, full interpretability of what agents have learned is essential.
\end{itemize}

\section{Project Significance}
\label{sec:project-significance}

This project contributes to both game AI and multi-agent reinforcement learning research:

\begin{itemize}
    \item \textbf{Emergent social dynamics:} Demonstrates that cooperation can arise purely from reward structure and learning, without being scripted --- a result relevant to computational social science and mechanism design.
    \item \textbf{Interpretable multi-agent RL:} By using tabular Q-learning (6,300 Q-values per NPC), every learned policy is directly inspectable. This contrasts with deep RL approaches where policy interpretation requires post-hoc explanation methods.
    \item \textbf{Decomposed reward framework:} The $R_i(t) = P_i(t) + G_i(t) + \lambda_i(t) \cdot C(t)$ formulation with dynamic $\lambda$ is domain-agnostic and transferable to industrial multi-agent coordination problems (factory automation, automotive ECU tuning).
    \item \textbf{Research playground:} The system provides a complete experimental platform with ablation controls, analytics dashboard, deterministic seeding, and automated bot playthroughs for reproducible research.
\end{itemize}

\section{Scope of Work}
\label{sec:scope}

The system is structured in three layers:

\begin{itemize}
    \item \textbf{Hierarchical MDP:} A 7-stage quest graph with 16 static checkpoints drives game progression. Dynamic checkpoints handle player deviations with nudging and convergence mechanisms.
    \item \textbf{Q-Learning NPC Agents:} Six NPCs run tabular Q-learning over a 225-state $\times$ 28-action space with role-specific action masking and adaptive personality coefficients.
    \item \textbf{Optional LLM Narration:} An HTTP adapter connects to OpenRouter, Ollama, or any OpenAI-compatible provider. Every LLM call has a deterministic template fallback, ensuring the game runs identically without a model.
\end{itemize}

The project is implemented as a web application (Python/FastAPI backend, vanilla JavaScript frontend) with a real-time analytics dashboard for monitoring cooperation curves, reward decomposition, and shock response.

\section{Objectives}
\label{sec:objectives}

The main objectives of this project are:

\begin{enumerate}
    \item Design and implement a multi-agent reinforcement learning environment within a medieval village RPG setting.
    \item Develop a decomposed reward model with a dynamic $\lambda$ coefficient that creates feedback loops between individual and community welfare.
    \item Implement six tabular Q-learning NPC agents with role-specific action masking and adaptive personality dynamics.
    \item Build a system shock engine to test agent resilience under non-stationary environmental perturbations.
    \item Conduct a systematic ablation study isolating the contribution of each subsystem to emergent cooperation.
    \item Integrate optional LLM narration with guardrails and template fallbacks for atmospheric game text.
\end{enumerate}

\section{Novelty}
\label{sec:novelty}

The key novel contributions of this project are:

\begin{itemize}
    \item \textbf{Dynamic $\lambda$ feedback loop:} Unlike static reward weighting in prior MARL work, the $\lambda_i(t) = 0.05 + \text{coop}_i \times 0.55$ mechanism creates a virtuous cycle where cooperation self-amplifies through the reward structure. Our ablation shows this contributes $+0.158$ to cooperation --- the second-largest factor after RL itself.
    \item \textbf{Fully interpretable multi-agent cooperation:} While prior game AI research uses deep RL (opaque policies), our tabular approach lets researchers directly inspect every Q-value and trace exactly why an NPC chose a cooperative action.
    \item \textbf{Integrated research playground:} The system combines RL agents, environmental shocks, adaptive personalities, quest progression, LLM narration, and a complete analytics dashboard in a single reproducible platform with ablation controls --- a combination not found in existing game AI testbeds.
\end{itemize}

% =============================================================================
% CHAPTER 2: LITERATURE SURVEY
% =============================================================================
\chapter{Literature Survey}
\label{chap:literature}

This chapter reviews prior work across four domains that converge in our project: multi-agent reinforcement learning, reward shaping, NPC intelligence in games, and LLM-enhanced narration. For each category, we summarize the major approaches and identify the gaps our project addresses.

\section{Multi-Agent Reinforcement Learning}
\label{sec:marl}

Cooperative MARL has been studied through several frameworks. QMIX~\cite{rashid2018qmix} factorizes joint Q-values under a monotonicity constraint for cooperative tasks; MADDPG~\cite{lowe2017maddpg} extends actor-critic methods with centralized critics during training; and Leibo et al.~\cite{leibo2017multi} demonstrate that cooperation or defection can emerge from independent learners depending on environmental incentives. Our approach uses independent Q-learners with a shared community reward signal and a dynamic $\lambda$ coefficient that gates how much each agent weights the community signal based on its own cooperation history.

\section{Reward Shaping}
\label{sec:reward-shaping}

Ng et al.~\cite{ng1999rewardshaping} proved that potential-based reward shaping preserves optimal policies while accelerating learning. Our decomposed reward model extends this idea: the community component $\lambda_i(t) \cdot C(t)$ acts as a shaping signal that nudges agents toward prosocial behavior, with the dynamic $\lambda$ controlling the strength of this nudge based on each agent's cooperation trajectory.

\section{NPC Intelligence in Games}
\label{sec:npc-intelligence}

The design of believable NPCs has evolved through several distinct paradigms over the past two decades. Traditional NPC behavior relies on behavior trees~\cite{isla2005bt} and hierarchical task networks, which require designers to manually author every decision branch. While effective for tightly controlled experiences, these approaches produce deterministic, repetitive behavior and scale poorly as game worlds grow in complexity. Bartle~\cite{bartle2004worlds} provides foundational principles for designing virtual worlds with believable inhabitants, emphasizing the importance of NPC autonomy and emergent social structures for player immersion --- goals that scripted systems struggle to achieve.

The shift toward learning-based approaches has opened new possibilities. Riedl and Bulitko~\cite{riedl2013interactive} survey intelligent systems approaches to interactive narrative, identifying the tension between author control and emergent behavior --- a central challenge our hierarchical MDP addresses by providing structured quest progression while allowing NPC behavior to emerge freely within each stage. More recently, Park et al.~\cite{park2023generative} demonstrated generative agents that simulate human behavior using LLM-driven memory and planning, producing remarkably believable social interactions in a sandbox environment. Their work shows that LLMs can serve as the cognitive backbone for NPCs, but at the cost of requiring persistent model access and incurring high inference costs.

Our system bridges the gap between these paradigms. Rather than replacing behavior trees entirely with learning (which sacrifices designer control) or using LLMs as the primary decision engine (which creates model dependency), we use tabular Q-learning for NPC decision-making within a structured quest framework, with LLMs providing an optional narration layer. This preserves full interpretability of NPC policies, maintains gameplay without model access, and still produces emergent cooperative behavior.

\section{LLM-Enhanced Game Narration}
\label{sec:llm-narration}

Large language models have been applied to interactive fiction and game content generation through several distinct approaches. Urbanek et al.~\cite{urbanek2019light} developed the LIGHT environment for grounded dialogue in a text-adventure setting, training models that generate dialogue conditioned on character personas and environmental descriptions. However, their system treats the language model as the primary interaction engine, making gameplay quality entirely dependent on model availability and output quality.

Ammanabrolu et al.~\cite{ammanabrolu2020story} address story realization --- expanding high-level plot events into coherent natural language sentences --- demonstrating that structured intermediate representations improve narrative coherence. This insight informs our template system: when the LLM is unavailable, structured templates serve as the ``plot event'' layer, ensuring that the game always communicates essential information. V\"{a}rtinen et al.~\cite{vartinen2022quests} explore GPT-based quest generation for role-playing games, showing that language models can produce contextually appropriate quest content but require careful prompting and validation to maintain game-mechanical consistency --- a finding that directly motivated our 9-stage guardrail pipeline.

Li et al.~\cite{li2016dialogue} apply deep reinforcement learning to dialogue generation, optimizing for long-term conversational quality rather than single-turn plausibility. Their work highlights the importance of grounding dialogue in context, a principle we apply through our contextual grounding metric that measures how well narration reflects current game state. Tambwekar et al.~\cite{tambwekar2019controllable} demonstrate controllable neural story generation via reward shaping, showing that narrative quality and plot adherence can be jointly optimized. Their approach of using reward signals to guide generation parallels our use of game-state context to constrain LLM output.

A key architectural decision in our system is the separation of LLM narration from core game logic. Unlike systems that use LLMs as the game engine (where a model failure breaks gameplay), our LLM serves as an optional enhancement layer. Every LLM code path has a deterministic template fallback, and a 9-stage guardrail pipeline validates all model outputs before they reach the player. This design ensures that the game runs identically with or without model access, while still benefiting from richer, more contextually grounded narration when a model is available.

\section{Comparison with Our Approach}
\label{sec:lit-comparison}

\begin{table}[H]
\centering
\caption{Comparison of Related Approaches and Our System}
\label{tab:lit-comparison}
\small
\begin{tabularx}{\textwidth}{|l|l|X|X|}
\hline
\textbf{Work} & \textbf{Year} & \textbf{Methodology} & \textbf{Our Differentiation} \\
\hline
QMIX~\cite{rashid2018qmix} & 2018 & Factorized joint Q-values, centralized training & We use independent learners with shared community reward --- no centralized critic needed \\
\hline
MADDPG~\cite{lowe2017maddpg} & 2017 & Centralized critic, decentralized actors (deep RL) & Tabular Q-learning for full interpretability; dynamic $\lambda$ instead of static reward mixing \\
\hline
Leibo et al.~\cite{leibo2017multi} & 2017 & Emergent social dilemmas in grid-worlds & Full RPG environment with quests, shocks, adaptive personalities, and LLM narration \\
\hline
Ng et al.~\cite{ng1999rewardshaping} & 1999 & Potential-based reward shaping theory & Extended with dynamic gating via $\lambda_i(t)$ that evolves with cooperation tendency \\
\hline
Isla~\cite{isla2005bt} & 2005 & Behavior trees for NPC AI & RL replaces hand-authored trees; cooperation emerges rather than being scripted \\
\hline
Park et al.~\cite{park2023generative} & 2023 & LLM-driven generative agents & LLM is optional narration layer, not cognitive backbone; full playability without model \\
\hline
LIGHT~\cite{urbanek2019light} & 2019 & LLM-driven interactive fiction & LLM is optional narration layer, not the game engine; 9-stage guardrails \\
\hline
Ammanabrolu et al.~\cite{ammanabrolu2020story} & 2020 & Story realization from plot events & Template fallbacks serve as structured plot-event layer; LLM enhances, not replaces \\
\hline
\end{tabularx}
\end{table}

% =============================================================================
% CHAPTER 3: METHODOLOGY
% =============================================================================
\chapter{Methodology}
\label{chap:methodology}

This chapter details the technical design of the system: the hierarchical MDP for quest progression, the Q-learning architecture for NPC agents, the decomposed reward model with dynamic $\lambda$, and the supporting subsystems (action masking, adaptive personalities, shock engine, LLM integration, and the turn processing pipeline).

\section{Hierarchical MDP}
\label{sec:hierarchical-mdp}

The quest is modeled as a two-level hierarchical MDP:

\textbf{Macro MDP (top level):}
\begin{equation}
    S_1 \to S_2 \to S_3 \to S_4 \to S_5 \to S_6 \to S_7 \to S_{\text{success}}
\end{equation}

Each macro stage $S_k$ contains a micro MDP $\mathcal{M}_{\text{micro}}^{(k)}$ with its own checkpoints, completion conditions, and transition rewards.

\textbf{Quest Stage Descriptions.} The 7 stages span the full narrative arc of a village quest:

\begin{enumerate}[nosep]
    \item \textbf{Arrival at the Gate ($S_1$):} The player enters the village through the main gate. Guards Aldric and Bryn are present and may interact with the player. The stage completes when the player moves past the gate into the village proper. This stage establishes the player's presence in the world and introduces the guard NPCs.
    \item \textbf{Seeking the Elder ($S_2$):} The player must navigate to Elder Maren's house and initiate conversation. Along the way, the player encounters other villagers and can engage in optional social interactions. The stage completes upon speaking with Elder Maren, who serves as the quest-giver.
    \item \textbf{Learning the Quest ($S_3$):} Elder Maren describes the village's problem and assigns the quest objective. The player receives context about what needs to be done and where to investigate. This is primarily a dialogue-driven stage where the LLM narration layer can produce the richest contextual descriptions.
    \item \textbf{Investigation ($S_4$):} The player explores the village to gather information relevant to the quest. This may involve talking to multiple NPCs, examining locations, and piecing together clues. NPC cooperation dynamics are most visible here, as the player observes how NPCs interact with each other and respond to inquiries.
    \item \textbf{Retrieval ($S_5$):} The player must obtain or accomplish the quest objective, which may involve navigating to a specific location, acquiring an item, or completing a task. Environmental shocks (famine, raids) can complicate this stage by altering NPC availability and willingness to assist.
    \item \textbf{Return ($S_6$):} The player returns to Elder Maren with the quest objective completed. This stage tests navigation and may involve additional NPC encounters along the return path.
    \item \textbf{Resolution ($S_7$):} Elder Maren acknowledges quest completion, and the village responds to the outcome. NPC adaptation states may shift based on the quest result, and community welfare is updated accordingly.
\end{enumerate}

\textbf{Dynamic Checkpoints.} When the player deviates from the expected path, the system generates dynamic checkpoints that acknowledge the deviation and guide convergence. The nudge mechanism works as follows: at each turn where the player is off the expected path, the system computes the distance to the nearest static checkpoint and applies a nudge reward:
\begin{equation}
    R_{\text{nudge}} = 0.3 \cdot (d_{\text{before}} - d_{\text{after}})
\end{equation}
This provides a gradient signal toward the nearest static checkpoint --- the player receives positive reward for moving closer and negative reward for moving further away. The coefficient 0.3 was tuned to provide noticeable guidance without overwhelming the intrinsic exploration rewards. If the player accumulates 5 deviations without returning to the expected path, forced convergence triggers: the system generates a bridge checkpoint that connects the player's current state to the nearest reachable static checkpoint, effectively re-routing the quest without breaking narrative coherence.

\textbf{Quest Auto-Restart.} When the quest completes before the turn limit, the quest state resets while preserving all NPC Q-tables and adaptation states, enabling multi-episode learning within a single session. This is critical for the ablation study: 500 turns typically spans 3--5 complete quest cycles, giving NPCs sufficient training episodes for cooperation to emerge. The auto-restart mechanism preserves the full NPC state (Q-tables, adaptation coefficients, personality parameters) while resetting only the quest-tracking variables (current stage, visited checkpoints, deviation count).

\section{NPC State Space and Q-Learning}
\label{sec:q-learning}

Each NPC maintains a tabular Q-table over a discretized state space:

\begin{table}[H]
\centering
\caption{NPC State Space Dimensions}
\label{tab:state_space}
\begin{tabular}{lll}
\toprule
\textbf{Dimension} & \textbf{Values} & \textbf{Count} \\
\midrule
Location & gate, village\_center, elders\_house, fields, tavern & 5 \\
Time & morning, midday, afternoon, evening, night & 5 \\
Energy & low ($<$33\% HP), medium (33--66\%), high ($>$66\%) & 3 \\
Mood & unhappy ($<$4), content (4--7), cheerful ($>$7) & 3 \\
\bottomrule
\end{tabular}
\end{table}

Total state space: $5 \times 5 \times 3 \times 3 = 225$ states. Action space: 28 universal actions across 7 categories (navigation, exploration, talk, social, combat, stealth, utility). Q-table size per NPC: 6,300 entries.

\textbf{Q-Learning update:}
\begin{equation}
    Q(s, a) \leftarrow Q(s, a) + \alpha \left[ r + \gamma \max_{a'} Q(s', a') - Q(s, a) \right]
\end{equation}
with $\alpha = 0.1$, $\gamma = 0.9$, $\varepsilon$-greedy exploration decaying from 0.15 to 0.05 at rate $0.995$/turn.

\section{Decomposed Reward Model}
\label{sec:reward-model}

Each NPC $i$ receives a reward signal composed of three terms:
\begin{equation}
    R_i(t) = P_i(t) + G_i(t) + \lambda_i(t) \cdot C(t)
    \label{eq:reward}
\end{equation}

\begin{itemize}[nosep]
    \item $P_i(t)$: \textbf{Penalty} --- $-10$ when HP drops below 20\%.
    \item $G_i(t)$: \textbf{Individual reward} --- $\sum_k w_k \cdot (\text{stat}_k^{\text{new}} - \text{stat}_k^{\text{old}})$, weighted by archetype-specific preferences.
    \item $C(t)$: \textbf{Community reward} --- Village-level welfare computed as $C(t) = 0.5 \cdot C_{\text{abs}}(t) + 0.5 \cdot C_{\Delta}(t)$, using centered non-linear shaping.
    \item $\lambda_i(t)$: \textbf{Dynamic lambda} --- $\lambda_i = 0.05 + \text{coop}_i \times 0.55$, ranging from 0.05 (selfish) to 0.60 (cooperative). Creates a feedback loop: positive community reward $\to$ cooperation rises $\to$ $\lambda$ increases $\to$ more weight on community reward.
\end{itemize}

\section{Role-Specific Action Masking}
\label{sec:action-masking}

Each NPC archetype has a preferred action set. Hard masking sets $Q(s,a) = -\infty$ for invalid actions, while soft masking applies a $+0.5$ Q-value bonus for role-aligned actions and a $-0.3$ penalty for misaligned ones. This provides structural inductive bias that improves learning efficiency without eliminating exploratory flexibility.

\section{Adaptive Personality Dynamics}
\label{sec:adaptation}

Four adaptation coefficients evolve per turn at base rate 0.018 with homeostatic drift rate 0.004 toward 0.5:

\begin{table}[H]
\centering
\caption{Adaptation Coefficients}
\label{tab:adaptation}
\begin{tabular}{lll}
\toprule
\textbf{Coefficient} & \textbf{Drives} & \textbf{Role Multipliers} \\
\midrule
cooperation\_tendency & $\lambda$ (community weight) & Gated by $\lambda$: $r_{\text{coop}} = r_{\text{base}} \times (0.5 + \lambda_i)$ \\
risk\_aversion & Risk avoidance & Guard 1.5$\times$, Farmer 1.2$\times$, Elder 0.6$\times$ \\
social\_sensitivity & Social feedback & Elder 1.4$\times$, Tavern 1.2$\times$, Guard 0.5$\times$ \\
shock\_resilience & Shock adaptation & Increases under shock pressure \\
\bottomrule
\end{tabular}
\end{table}

\section{System Shock Engine}
\label{sec:shock-engine}

Five shock types inject non-stationary perturbations:

\begin{table}[H]
\centering
\caption{System Shock Catalog}
\label{tab:shocks}
\begin{tabular}{llll}
\toprule
\textbf{Shock} & \textbf{Duration} & \textbf{Reward Scale} & \textbf{Purpose} \\
\midrule
Famine & 15 turns & 0.5$\times$ & Resource scarcity \\
Bandit Raid & 8 turns & 0.6$\times$ & External threat \\
Plague & 12 turns & 0.4$\times$ & Sustained decline \\
Trade Boom & 10 turns & 1.5$\times$ & Recovery (positive) \\
Harsh Winter & 20 turns & 0.7$\times$ & Long-duration stress \\
\bottomrule
\end{tabular}
\end{table}

Shocks directly pressure cooperation downward: when shock intensity exceeds 0.1, cooperation tendency decreases proportionally, independent of the reward signal.

\section{LLM Integration}
\label{sec:llm-integration}

The LLM integration layer is designed as a thin, provider-agnostic HTTP adapter that cleanly separates game logic from narration generation. This section details the architecture, use cases, guardrail pipeline, fallback system, and model selection rationale.

\subsection{HTTP Adapter Architecture}

The LLM service abstracts over three provider backends through a unified HTTP adapter:

\begin{itemize}[nosep]
    \item \textbf{Ollama:} Local inference via the Ollama REST API (\texttt{/api/generate}). Used for development and offline play.
    \item \textbf{OpenRouter:} Cloud inference through the OpenRouter aggregation service, which proxies to multiple model providers. Enables access to models not available locally.
    \item \textbf{LM Studio / OpenAI-compatible:} Any server exposing the OpenAI \texttt{/v1/chat/completions} endpoint. This includes LM Studio, vLLM, and text-generation-inference.
\end{itemize}

The adapter normalizes all provider-specific details (authentication headers, request format, response parsing) behind a single \texttt{generate()} interface. The calling code never knows which provider is active --- it submits a prompt and receives text. Provider selection is configured at startup via environment variables, and the system gracefully falls back to templates if no provider is reachable.

\subsection{Use Cases}

The LLM is invoked in four distinct contexts within the game:

\begin{enumerate}[nosep]
    \item \textbf{Input Parsing:} The player types free-form natural language commands. The LLM parses these into structured action objects (action type, target, parameters). When the LLM is unavailable, a spaCy-based NLP parser provides the fallback. This is the most latency-sensitive use case, as the player waits for the parse result before seeing the game response.
    \item \textbf{NPC Dialogue:} When the player talks to an NPC, the LLM generates contextual dialogue based on the NPC's role, personality state, current location, time of day, and recent events. The prompt includes the NPC's adaptation coefficients so that a highly cooperative NPC speaks differently from a newly encountered one. The template fallback produces role-appropriate but generic dialogue lines.
    \item \textbf{Checkpoint Generation:} Dynamic checkpoints (generated when the player deviates from the expected quest path) use the LLM to produce narratively coherent bridge text that acknowledges the deviation and guides the player back. The template fallback generates functional but less atmospheric checkpoint descriptions.
    \item \textbf{Narration Enhancement:} The primary use case. After each turn is resolved by the game engine, the LLM receives the raw game event data (action taken, location, NPCs present, time, weather, active shocks) and produces atmospheric narration. This is the most forgiving use case for latency, as narration can be generated asynchronously while the UI displays the mechanical result.
\end{enumerate}

\subsection{9-Stage Guardrail Pipeline}

Every LLM output passes through a 9-stage validation pipeline before reaching the player:

\begin{enumerate}[nosep]
    \item \textbf{Strip Reasoning Blocks:} Some models (particularly those with chain-of-thought training) emit \texttt{<think>...</think>} or reasoning preambles. This stage removes any such blocks, extracting only the intended output.
    \item \textbf{JSON Parsing:} For structured outputs (input parsing, checkpoint generation), the raw text is parsed as JSON. Malformed JSON triggers an immediate fallback to templates.
    \item \textbf{Schema Validation:} The parsed JSON is validated against the expected schema for the specific use case (e.g., action objects must have \texttt{type}, \texttt{target}, and \texttt{parameters} fields). Missing or extra fields trigger fallback.
    \item \textbf{Enum Validation:} String values are checked against valid enumerations. For example, action types must be one of the 28 defined actions; locations must be one of the 5 village areas. Invalid values trigger fallback.
    \item \textbf{Value Clamping:} Numeric values are clamped to valid ranges. Stat modifiers, reward values, and other numeric outputs are bounded to prevent the LLM from generating extreme values that would break game balance.
    \item \textbf{Action Reference Validation:} For outputs that reference game entities (NPCs, items, locations), the pipeline verifies that all referenced entities actually exist in the current game state. A narration that mentions an NPC who is not present at the current location is rejected.
    \item \textbf{Content Filtering:} Output is scanned for inappropriate content, out-of-setting language (modern slang, real-world references), and narrative inconsistencies. Flagged content triggers template fallback.
    \item \textbf{Length Enforcement:} Output is truncated or rejected if it exceeds maximum length limits. Narrations are capped to prevent the LLM from generating excessively verbose text that would overwhelm the UI.
    \item \textbf{Final Sanitization:} Remaining edge cases are handled: encoding normalization, whitespace cleanup, and escape character removal. The output is ready for display.
\end{enumerate}

If any stage fails, the pipeline short-circuits and returns the deterministic template output for that specific use case. This fail-fast design ensures that a single malformed LLM response never propagates to the player.

\subsection{Template Fallback System}

Every code path that calls the LLM has a co-located template fallback. Templates are parameterized strings that incorporate game-state variables (location name, NPC name, time of day, action type) to produce contextually relevant --- though formulaic --- text. The fallback system guarantees that the game is fully playable, narratively coherent, and mechanically identical regardless of whether an LLM is available. This separation is a deliberate architectural choice: the game engine produces correct gameplay; the LLM produces richer narration. Removing the LLM degrades prose quality but never gameplay quality.

\subsection{Model Selection}

Two models were selected for their balance of quality, speed, and resource requirements:

\begin{itemize}[nosep]
    \item \textbf{Local inference:} \texttt{qwen3:4b} via LM Studio. This 4B-parameter model runs on consumer hardware (8GB VRAM) with acceptable latency ($<$2 seconds per generation). Despite its small size, the guardrail pipeline ensures output quality remains high --- the 9 validation stages compensate for the model's occasional inconsistencies.
    \item \textbf{Cloud inference:} \texttt{meta-llama/llama-3.2-3b-instruct:free} via OpenRouter. This provides a zero-cost cloud option for users without local GPU resources. The 3B-parameter model produces slightly less varied narration than \texttt{qwen3:4b} but passes the guardrail pipeline at comparable rates.
\end{itemize}

The choice of small models (3--4B parameters) is deliberate: larger models produce marginally better prose but at significantly higher latency and cost, which undermines the real-time game experience. The guardrail pipeline allows small models to achieve narration quality sufficient for immersive gameplay.

\section{Turn Processing Pipeline}
\label{sec:turn-pipeline}

Each player action triggers a 10-step pipeline that orchestrates the complete game loop:
\begin{enumerate}[nosep]
    \item Resolve player action (combat, social, navigation, etc.)
    \item Log event with importance scoring
    \item Check quest completion or deviation
    \item Process 6 NPC turns (Q-learning action selection, reward computation, Q-table update)
    \item Check random events
    \item Shock engine tick (decay intensities, apply stat drain, trust modifiers)
    \item Time advancement (4 turns = 1 period)
    \item Reputation decay, stamina regeneration, perception checks
    \item Game-over conditions (or quest restart if turns remain)
    \item Telemetry logging (world snapshot, reward decomposition, adaptation state)
\end{enumerate}

The data flow through the pipeline is strictly sequential: each step depends on the state produced by the previous step. Step~1 resolves the player's action and modifies the world state (player position, NPC stats, inventory). Step~2 logs this event with an importance score used by the LLM narration system to decide how much detail to include. Step~3 checks whether the player action advanced or deviated from the quest, triggering dynamic checkpoint generation if needed. Step~4 is the core RL step: each of the 6 NPCs observes its current state, selects an action via $\varepsilon$-greedy over its Q-table (with role masking applied), executes the action, receives the decomposed reward $R_i(t)$, and updates its Q-table. Step~5 checks for random environmental events (merchant arrivals, weather changes). Step~6 advances the shock engine: active shocks decay in intensity, stat drains are applied to affected NPCs, and trust modifiers adjust NPC-player relationships. Step~7 advances the in-game clock (every 4 turns constitutes one time period transition, e.g., morning to midday). Step~8 applies passive effects: reputation naturally decays toward neutral, stamina regenerates, and NPCs run perception checks to notice nearby events. Step~9 evaluates terminal conditions (player death, quest completion with no remaining turns) or triggers quest auto-restart if turns remain. Finally, Step~10 captures the complete world state as a telemetry snapshot for the analytics dashboard, recording reward decomposition, adaptation coefficients, and cooperation indices.

This pipeline ensures that every turn produces a consistent, reproducible sequence of state transitions. The deterministic ordering is critical for the ablation study: with the same seed and player strategy, the same sequence of turns produces identical results across runs.

% =============================================================================
% CHAPTER 4: CHALLENGES AND SOLUTIONS
% =============================================================================
\chapter{Challenges and Solutions}
\label{chap:challenges}

Developing a multi-agent RL system within a full game environment introduced several technical and design challenges. Table~\ref{tab:challenges} summarizes the major issues faced and their corresponding solutions.

\begin{table}[H]
\centering
\caption{Major Implementation Challenges and Their Solutions}
\label{tab:challenges}
\begin{tabularx}{\textwidth}{|X|X|}
\hline
\textbf{Challenge} & \textbf{Solution} \\
\hline
\textbf{Multi-agent non-stationarity:} Six Q-learners in a shared environment violate the stationary MDP assumption needed for convergence guarantees.
& Slow adaptation rate (0.018), homeostatic drift toward 0.5, and a shared community reward signal that aligns agents, empirically producing stable cooperation. \\
\hline
\textbf{Reward hacking risk:} Agents could exploit $C(t)$ via degenerate strategies that inflate village stats without genuine cooperation.
& Competing objectives (role-specific $G_i$, penalty $P_i$) and non-linear shaping (sigmoid, sub-linear) prevent saturation exploits. Ablation verified cooperation correlates with prosocial actions. \\
\hline
\textbf{LLM output quality:} Local 4B models produce inconsistent JSON, hallucinate game state, and occasionally generate inappropriate content.
& 9-stage guardrail pipeline (JSON parse, schema/enum validation, value clamping, content filtering) with deterministic template fallbacks on every code path. \\
\hline
\textbf{Reproducibility:} Python's hash randomization and platform-specific RNG behavior make cross-run reproducibility difficult.
& Single \texttt{MASTER\_SEED}; NPC pre-training seeds derived from registry index (not \texttt{hash()}). Explicit reproducibility tests verify identical playthroughs. \\
\hline
\textbf{Windows asyncio deadlocks:} Module-based execution caused \texttt{ProactorEventLoop} deadlocks with \texttt{asyncio.to\_thread}.
& Direct script execution with \texttt{sys.path} insertion, avoiding the Windows-specific event loop interaction. \\
\hline
\end{tabularx}
\end{table}

% =============================================================================
% CHAPTER 5: RESULTS AND DISCUSSION
% =============================================================================
\chapter{Results and Discussion}
\label{chap:results}

This chapter presents the experimental design, ablation study results, reward analysis, NPC adaptation patterns, and narrative quality evaluation.

\section{Experimental Design}
\label{sec:experimental-design}

\subsection{Ablation Conditions}

We define seven ablation conditions, each isolating one subsystem:

\begin{table}[H]
\centering
\caption{Ablation Conditions}
\label{tab:conditions}
\begin{tabular}{llccccc}
\toprule
\textbf{ID} & \textbf{Description} & \textbf{MDP} & \textbf{RL} & \textbf{Shocks} & \textbf{Mask} & \textbf{Dyn.\ $\lambda$} \\
\midrule
C1 & Full System (baseline) & \checkmark & \checkmark & \checkmark & \checkmark & \checkmark \\
C2 & No LLM & \checkmark & \checkmark & \checkmark & \checkmark & \checkmark \\
C3 & No RL (schedules only) & \checkmark & --- & \checkmark & --- & --- \\
C4 & Flat MDP (no hierarchy) & --- & \checkmark & \checkmark & \checkmark & \checkmark \\
C5 & No shocks & \checkmark & \checkmark & --- & \checkmark & \checkmark \\
C6 & No role masking & \checkmark & \checkmark & \checkmark & --- & \checkmark \\
C7 & Static $\lambda = 0.325$ & \checkmark & \checkmark & \checkmark & \checkmark & --- \\
\bottomrule
\end{tabular}
\end{table}

\subsection{Metrics}

\begin{itemize}[nosep]
    \item \textbf{Cooperation index} $\mathcal{K}(t) = \frac{1}{N} \sum_{i=1}^{N} \text{coop}_i(t)$ --- average cooperation tendency across all NPCs at each turn.
    \item \textbf{Social welfare index} --- average cooperation tendency weighted by community reward performance.
    \item \textbf{Policy entropy} $\mathcal{D}_i = -\sum_{a \in \mathcal{A}} p_i(a) \log p_i(a)$ --- computed from full per-action counts.
    \item \textbf{Contextual grounding} --- fraction of game-state keywords (action, location, NPC names) appearing in each narration.
    \item \textbf{Vocabulary diversity} --- type-token ratio (unique words / total words) across narrations.
    \item \textbf{Shock recovery time} --- turns after shock expiry until $\mathcal{K}(t)$ returns to pre-shock baseline.
\end{itemize}

\subsection{Protocol}

Automated bot playthroughs using the \texttt{quest\_focused} strategy with seed 42 for reproducibility. Each condition runs for 500 turns with quest auto-restart. A famine shock is injected at turn 166 for conditions with shocks enabled. Normal difficulty preset with adaptive difficulty disabled.

\section{Ablation Study Results}
\label{sec:ablation-results}

\begin{table}[H]
\centering
\caption{Ablation Study Results (500 turns, seed=42)}
\label{tab:results}
\begin{tabular}{llcc}
\toprule
\textbf{Condition} & \textbf{Description} & \textbf{Final $\mathcal{K}$} & \textbf{$\Delta$ vs C1} \\
\midrule
C1 & Full System (baseline) & 0.851 & --- \\
C3 & No RL (schedules only) & 0.500 & $-0.351$ \\
C4 & Flat MDP (no hierarchy) & 0.953 & $+0.102$ \\
C5 & No Shocks & 0.873 & $+0.022$ \\
C6 & No Role Masking & 0.763 & $-0.088$ \\
C7 & Static $\lambda = 0.325$ & 0.693 & $-0.158$ \\
\bottomrule
\end{tabular}
\end{table}

\begin{figure}[H]
    \centering
    \includegraphics[width=\textwidth]{images/fig_cooperation.png}
    \caption{Cooperation index $\mathcal{K}(t)$ over 500 turns for six ablation conditions. C3 (No RL) remains flat at 0.500 while all RL-enabled conditions show emergence. C7 (Static $\lambda$) rises slower than C1, confirming the feedback loop's amplifying effect. C4 (Flat MDP) rises fastest due to more frequent quest cycling.}
    \label{fig:cooperation}
\end{figure}

\subsection{Finding 1: RL Drives Cooperation}

C1 (Full System) achieves a cooperation index of 0.851 while C3 (No RL) remains at exactly 0.500 --- the initialization baseline. The 0.351 gap demonstrates that cooperation is an \textit{emergent} property of the learning process, not a byproduct of game mechanics.

\subsection{Finding 2: Role Masking Aids Learning}

C6 (No Role Masking) achieves 0.763, compared to C1's 0.851. Without role constraints, NPCs explore the full 28-action space uniformly, producing noisier Q-value estimates and slower convergence.

\subsection{Finding 3: Shock Resilience}

C5 (No Shocks) reaches 0.873, slightly above C1's 0.851. The famine shock at turn 166 creates a visible dip in C1's cooperation curve (\cref{fig:cooperation}), but NPCs recover within approximately 30 turns, demonstrating adaptive resilience.

\subsection{Finding 4: MDP Structure}

C4 (Flat MDP) achieves the highest cooperation at 0.953. Without hierarchical stage gates, quest completion is faster, giving NPCs more learning episodes in 500 turns. This suggests that quest structure affects the \textit{rate} of learning more than the \textit{direction}.

\subsection{Finding 5: Dynamic Lambda Amplifies Cooperation}

C7 (Static $\lambda = 0.325$) achieves 0.693, compared to C1's 0.851 --- a 0.158 difference, making the dynamic $\lambda$ feedback loop the second-largest contributor after RL itself. The virtuous cycle --- positive community outcomes $\to$ cooperation rises $\to$ $\lambda$ increases $\to$ more weight on community reward --- provides substantial acceleration.

\begin{figure}[H]
    \centering
    \includegraphics[width=0.85\textwidth]{images/fig_final_bar.png}
    \caption{Final cooperation index by condition after 500 turns. C7 (Static $\lambda$) falls between C6 and C3, confirming the feedback loop's amplifying role.}
    \label{fig:final_bar}
\end{figure}

\section{Reward Decomposition}
\label{sec:reward-decomposition}

\begin{figure}[H]
    \centering
    \includegraphics[width=\textwidth]{images/fig_reward.png}
    \caption{Average reward components over time for C1 (Full System). Community reward $C(t)$ dominates the total signal as $\lambda$ increases. The famine shock at turn 166 is visible as a dip in community reward.}
    \label{fig:reward}
\end{figure}

The reward decomposition confirms the hypothesized dynamic: early turns are dominated by individual reward $G(t)$; as stats stabilize, community reward $C(t)$ becomes the primary gradient signal, amplified by the rising $\lambda$ coefficient.

\section{NPC Adaptation and Role Divergence}
\label{sec:npc-adaptation}

\begin{figure}[H]
    \centering
    \includegraphics[width=\textwidth]{images/fig_npc_adaptation.png}
    \caption{NPC adaptation states after 500 turns. Left: C1 (Full System) shows role-specific divergence. Right: C3 (No RL) remains at baseline 0.500 across all coefficients.}
    \label{fig:npc_adaptation}
\end{figure}

\begin{table}[H]
\centering
\caption{NPC Final Adaptation States (C1, Turn 500)}
\label{tab:npc_states}
\begin{tabular}{llccccc}
\toprule
\textbf{NPC} & \textbf{Role} & \textbf{Coop.} & \textbf{Risk} & \textbf{Social} & \textbf{Resil.} & $\boldsymbol{\lambda}$ \\
\midrule
Elder Maren & Elder & 0.850 & 0.500 & 0.898 & 0.499 & 0.518 \\
Farmer Jak & Farmer & 0.851 & 0.494 & 0.898 & 0.499 & 0.518 \\
Aldric & Guard & 0.851 & 0.500 & 0.785 & 0.499 & 0.518 \\
Bryn & Guard & 0.851 & 0.495 & 0.785 & 0.499 & 0.518 \\
Tessa & Tavern Keeper & 0.851 & 0.495 & 0.898 & 0.499 & 0.518 \\
Old Petra & Villager & 0.851 & 0.484 & 0.898 & 0.499 & 0.518 \\
\bottomrule
\end{tabular}
\end{table}

The most notable pattern is cooperation convergence: all NPCs reach $\sim$0.851 regardless of role, indicating that the community reward provides a consistent gradient. Social sensitivity diverges by role as expected (Elder and Tavern Keeper at 0.898 vs Guards at 0.785), driven by role-specific multipliers.

\section{Action Distribution Shift}
\label{sec:action-shift}

\begin{figure}[H]
    \centering
    \includegraphics[width=0.85\textwidth]{images/fig_action_shift.png}
    \caption{Change in NPC action proportions between early and late game (C1). Green = actions used more in late game; red = used less. NPCs shift from social actions toward utility and exploration.}
    \label{fig:action_shift}
\end{figure}

NPCs shift from social actions (talk, trade) toward utility and exploration (examine, work, rest) as they learn which actions maximize their decomposed reward.

\section{Narrative Quality}
\label{sec:narrative-quality}

\begin{table}[H]
\centering
\caption{Narrative Quality: LLM (qwen3:4b) vs Templates (20 turns)}
\label{tab:narrative}
\begin{tabular}{lccc}
\toprule
\textbf{Metric} & \textbf{C1 (LLM)} & \textbf{C2 (Template)} & \textbf{Winner} \\
\midrule
Contextual grounding & 0.567 & 0.488 & LLM \\
Vocabulary diversity (TTR) & 0.323 & 0.277 & LLM \\
Unique words & 396 & 151 & LLM \\
Avg.\ narration length & 61.4 words & 27.3 words & LLM \\
Lexical coherence & 0.202 & 0.213 & Template \\
\bottomrule
\end{tabular}
\end{table}

\textbf{Example narration comparison:}

\begin{quote}
\textbf{LLM (C1):} ``You scan the Gate in the morning light, where shadows are still thin and the air hums with the quiet buzz of dawn, but no shelter offers true cover.''

\textbf{Template (C2):} ``You look for a hiding spot but can't find adequate cover. Aldric and Bryn observe your actions.''
\end{quote}

\subsection{Contextual Grounding and Immersion}

The LLM achieves a contextual grounding score of 0.567 compared to templates' 0.488 --- a 16\% improvement. Contextual grounding measures the fraction of game-state keywords (current location, time of day, NPCs present, active action, weather conditions) that appear in the narration. Higher grounding means the narration reflects what is actually happening in the game world rather than producing generic text. This matters for player immersion: when narration references the specific gate at dawn with guards nearby, the player feels that the game world is coherent and responsive. Templates achieve reasonable grounding because they are parameterized with game-state variables, but they cannot weave these elements into natural prose the way a language model can.

\subsection{Vocabulary Diversity vs Lexical Coherence}

The LLM produces 2.6$\times$ more unique words (396 vs 151) and achieves a higher type-token ratio (0.323 vs 0.277), indicating significantly richer vocabulary. However, templates score marginally higher on lexical coherence (0.213 vs 0.202). This trade-off is expected: templates reuse identical phrases across narrations, creating artificial token overlap that inflates coherence scores. The LLM's lower coherence score actually reflects its strength --- varied phrasing that avoids repetition. In practice, lexical variety produces more engaging prose: players notice and are bored by repeated phrases but appreciate diverse descriptions of similar events.

\subsection{Guardrail Pipeline and Small Model Quality}

A key finding is that the 9-stage guardrail pipeline enables small (4B-parameter) models to produce narration of sufficient quality for immersive gameplay. Without guardrails, \texttt{qwen3:4b} occasionally produces off-setting language, references nonexistent game entities, or generates excessively long responses. The pipeline catches these issues before they reach the player, falling back to templates for the specific turn where the model output failed validation. Over a 20-turn evaluation window, the fallback rate was below 15\%, meaning the LLM successfully contributed narration for the majority of turns. This demonstrates that the combination of a small local model with a robust validation pipeline is a viable alternative to using larger, more expensive models.

\subsection{Qualitative Observations}

Beyond the quantitative metrics, several qualitative differences are notable. LLM narrations tend to adopt a literary style with sensory details (light, sound, temperature) that templates omit entirely. The LLM occasionally produces narration that references the passage of time within the game world (``the shadows grow longer as midday approaches''), creating a sense of temporal progression. Templates, by contrast, describe only the immediate mechanical outcome of each action. When system shocks are active, the LLM adapts its tone --- famine produces narration with imagery of scarcity and hunger, while trade booms produce descriptions of bustling activity. Templates do not adjust tone for shocks, producing the same text regardless of environmental conditions.

\section{Social Welfare}
\label{sec:social-welfare}

\begin{figure}[H]
    \centering
    \includegraphics[width=\textwidth]{images/fig_welfare.png}
    \caption{Village social welfare index over 500 turns. Tracks cooperation emergence at the village level, weighted by community reward performance.}
    \label{fig:welfare}
\end{figure}

% =============================================================================
% CHAPTER 6: TOOLS, TECHNOLOGIES, AND PARAMETERS
% =============================================================================
\chapter{Tools, Technologies, and Parameters}
\label{chap:tools}

This chapter documents the technology stack, NPC configuration, and key hyperparameters used in the implementation. The system comprises approximately 17,700 lines of Python (backend) and 6,300 lines of JavaScript (frontend).

\section{Tools and Technologies}
\label{sec:tools-tech}

\begin{itemize}
    \item \textbf{Programming Languages / Frameworks:} Python 3.12 (backend), Vanilla JavaScript (frontend), FastAPI (web framework), Pydantic (data validation).
    \item \textbf{RL Engine:} NumPy for tabular Q-table operations.
    \item \textbf{NLP:} spaCy for input parsing fallback.
    \item \textbf{LLM Integration:} HTTP adapter supporting Ollama, OpenRouter, and LM Studio (OpenAI-compatible endpoints).
    \item \textbf{Frontend Visualization:} Chart.js (analytics dashboard), Cytoscape.js (real-time MDP graph).
    \item \textbf{Research Analytics:} Matplotlib for experiment figures and ablation plots.
    \item \textbf{Testing:} pytest (91 tests across 5 suites: core gameplay, shock engine, adaptation \& masking, analytics, integration).
    \item \textbf{Package Management:} \texttt{uv} with \texttt{pyproject.toml}.
\end{itemize}

\section{NPC Roster and Configuration}
\label{sec:npc-roster}

\begin{table}[H]
\centering
\caption{NPC Roster and Reward Weights}
\label{tab:npc_roster}
\begin{tabular}{llllllll}
\toprule
\textbf{NPC} & \textbf{Role} & \textbf{Location} & \textbf{Happiness} & \textbf{Income} & \textbf{Health} & \textbf{Rep.} \\
\midrule
Elder Maren & Elder & Elder's House & 0.3 & 0.1 & 0.2 & 0.4 \\
Farmer Jak & Farmer & Fields & 0.2 & 0.4 & 0.2 & 0.2 \\
Aldric & Guard & Gate & 0.1 & 0.2 & 0.3 & 0.4 \\
Bryn & Guard & Gate & 0.1 & 0.2 & 0.3 & 0.4 \\
Tessa & Tavern Keeper & Tavern & 0.2 & 0.4 & 0.1 & 0.3 \\
Old Petra & Villager & Village Center & 0.4 & 0.2 & 0.2 & 0.2 \\
\bottomrule
\end{tabular}
\end{table}

\section{Key Parameters and Settings}
\label{sec:parameters}

\begin{table}[H]
\centering
\caption{Important Parameters Used in Implementation}
\label{tab:parameters}
\begin{tabularx}{\textwidth}{|l|X|}
\hline
\textbf{Parameter} & \textbf{Value / Description} \\
\hline
Learning rate ($\alpha$) & 0.1 --- standard tabular Q-learning rate \\
\hline
Discount factor ($\gamma$) & 0.9 --- balances immediate vs future reward \\
\hline
Exploration ($\varepsilon$) & Decays from 0.15 to 0.05 at rate 0.995/turn \\
\hline
Pre-training & 100 episodes $\times$ 50 turns per NPC at $\varepsilon = 0.5$ \\
\hline
Adaptation rate & 0.018 base rate with homeostatic drift 0.004 toward 0.5 \\
\hline
Dynamic $\lambda$ & $\lambda_i = 0.05 + \text{coop}_i \times 0.55$, range [0.05, 0.60] \\
\hline
Static $\lambda$ (C7) & Fixed at 0.325 (midpoint of dynamic range) \\
\hline
Soft mask bonus/penalty & $+0.5$ for role-aligned actions, $-0.3$ for misaligned \\
\hline
HP penalty threshold & $-10$ reward when HP drops below 20\% \\
\hline
Nudge reward & $R_{\text{nudge}} = 0.3 \cdot (d_{\text{before}} - d_{\text{after}})$ \\
\hline
Master seed & 42 (all randomness derived deterministically) \\
\hline
Experiment duration & 500 turns per condition with quest auto-restart \\
\hline
Famine shock injection & Turn 166 (one-third of session) \\
\hline
\end{tabularx}
\end{table}

% =============================================================================
% CHAPTER 7: CONCLUSION AND FUTURE SCOPE
% =============================================================================
\chapter{Conclusion and Future Scope}
\label{chap:conclusion}

This project demonstrates that emergent cooperative behavior can arise from independent Q-learning agents when the reward signal includes a community welfare component gated by a dynamic $\lambda$ coefficient. The key findings, with emphasis on game architecture and narration insights:

\begin{enumerate}
    \item \textbf{The hierarchical quest structure provides meaningful game progression while enabling RL.} The 7-stage MDP with dynamic checkpoints creates a coherent narrative arc that guides player experience, while quest auto-restart ensures NPCs receive sufficient training episodes. The C4 result ($\mathcal{K} = 0.953$ for flat MDP) shows that quest structure affects the \textit{rate} of cooperation emergence but not its \textit{direction} --- a useful property for game designers who need both narrative structure and adaptive NPCs.
    \item \textbf{LLM narration adds quality without sacrificing reliability.} The optional narration layer achieves higher contextual grounding (0.57 vs 0.49) and 2.6$\times$ richer vocabulary than templates, while the 9-stage guardrail pipeline ensures that small 4B-parameter models produce output of sufficient quality. The template fallback system guarantees identical gameplay with or without model access.
    \item \textbf{Cooperation emerges from learning.} $\mathcal{K}(t)$ rises from 0.500 to 0.851 over 500 turns with Q-learning, while schedule-only NPCs show no change ($+0.351$ gap). This confirms the research hypothesis.
    \item \textbf{The dynamic $\lambda$ feedback loop is a primary amplifier.} Static $\lambda$ still produces emergence (0.500 $\to$ 0.693), but dynamic $\lambda$ adds $+0.158$ through the virtuous cycle --- the second-largest single-component contribution after RL itself.
    \item \textbf{The turn processing pipeline produces reproducible, deterministic game loops.} The 10-step pipeline with telemetry logging enables the ablation study methodology: identical seeds produce identical playthroughs, allowing clean isolation of each subsystem's contribution.
\end{enumerate}

\section{Limitations}
\label{sec:limitations}

\begin{itemize}[nosep]
    \item Results are from automated bot playthroughs with a single strategy (\texttt{quest\_focused}); human player behavior may produce different cooperation dynamics.
    \item The adaptation rate (0.018) and role-specific multipliers were empirically tuned; systematic sensitivity analysis is needed.
    \item Tabular Q-learning limits scalability --- 225 states is manageable for 6 NPCs but would not scale to larger worlds without function approximation.
    \item Narrative coherence metrics rely on token overlap rather than semantic similarity; sentence-transformer embeddings would provide stronger signal.
\end{itemize}

\section{Future Scope}
\label{sec:future-scope}

\begin{itemize}[nosep]
    \item Scale to 100 bot playthroughs per condition with varied strategies for statistical significance.
    \item Conduct user studies (20 participants per condition) with post-session questionnaires.
    \item Compare tabular Q-learning vs.\ DQN for NPC behavior.
    \item Extend to branching quest graphs with multiple concurrent quest lines.
    \item Explore longer sessions ($>$1000 turns) to test cooperation saturation.
\end{itemize}

\section{Real-World Transfer}
\label{sec:real-world-transfer}

The decomposed reward framework $R_i(t) = P_i(t) + G_i(t) + \lambda_i(t) \cdot C(t)$ generalises beyond games to any domain where heterogeneous subsystems must balance local objectives against a shared global target.

\textbf{Smart Factory Coordination.}
In an automated production line, each cell (stamping, welding, painting) optimises a local KPI---throughput, defect rate, idle time---analogous to role-specific $G_i$ weights. The community signal $C(t)$ maps to plant-wide OEE (Overall Equipment Effectiveness). Supply-chain disruptions, equipment failures, and demand spikes correspond directly to famine, plague, and trade-boom shocks. Current Industry~4.0 approaches rely on centralised MPC controllers that must be re-modelled when the line changes; a dynamic-$\lambda$ architecture lets each cell run a lightweight learner and join the reward structure without replanning the plant. Our C1~vs.~C3 result ($+0.351$ cooperation from RL over fixed schedules) suggests learned policies would significantly outperform static dispatch rules.

\textbf{Automotive ECU Self-Tuning.}
A modern engine's sensor--actuator subsystems (knock sensor, O$_2$ sensor, throttle body, MAP, coolant temperature, crank position) map onto heterogeneous NPC roles: each maximises a local objective (detonation margin, fuel-trim accuracy, airflow response) while the community target $C(t)$ is engine-level efficiency (e.g.\ brake-specific fuel consumption). The dynamic $\lambda$ feedback loop enables continuous self-calibration: a factory base map provides the pre-trained Q-table, and on-road experience adjusts cooperation tendencies as altitude, fuel quality, and component wear shift the reward landscape. Component aging becomes a slow-onset shock; aftermarket modifications (larger turbo, free-flowing exhaust) reset part of the state, and the system re-converges without manual tuning---our C4 result ($\mathcal{K} = 0.953$) suggests rapid re-convergence is achievable.

\textbf{Bridging the Gap.}
Two challenges must be addressed: (i)~the 225-state tabular formulation must give way to function approximation (DQN or policy gradients) for high-dimensional, real-time sensor streams; and (ii)~the $P_i$ penalty discourages but does not formally prevent catastrophic states---a hybrid architecture enforcing hard constraints from classical control theory while allowing RL-based optimisation within the safe envelope is the likely path forward.

% =============================================================================
% BIBLIOGRAPHY
% =============================================================================
\addcontentsline{toc}{chapter}{Bibliography}
\bibliographystyle{plainnat}
\bibliography{references}

\end{document}
